% !TEX program = xelatex
% !TEX options = -synctex=1 -interaction=nonstopmode -file-line-error
% example-teaching.tex
% Demonstrates rossbeamerteaching features: learning objectives, poll question,
% worked example, key terms, exam question, recap.
\documentclass[palette=lecture]{rossbeamerteaching}

\title{Introduction to Synaptic Transmission}
\subtitle{NEUR2001 --- Foundations of Neuroscience}
\author{Dr.~James O'Brien}
\institute{School of Biological Sciences, University of Manchester}
\date{Week 4, Semester 1}
\module{NEUR2001: Foundations of Neuroscience}

\begin{document}

% --- Title ---
\begin{frame}[plain]
  \titlepage
\end{frame}

% --- Learning Objectives ---
\begin{frame}{Today's Session}
  \begin{learningobjectives}
    \item Describe the sequence of events during chemical synaptic transmission
    \item Distinguish between ionotropic and metabotropic receptors
    \item Explain the role of calcium in neurotransmitter release
    \item Identify the major excitatory and inhibitory neurotransmitters
  \end{learningobjectives}
\end{frame}

% --- Key Terms ---
\begin{keytermsframe}
  \keytermentry{Synapse}{The junction between two neurons where signals
    are transmitted via chemical or electrical means}
  \keytermentry{Vesicle}{Membrane-bound organelle containing
    neurotransmitter molecules}
  \keytermentry{SNARE complex}{Protein machinery mediating vesicle
    fusion with the presynaptic membrane}
  \keytermentry{Reuptake}{Process by which released neurotransmitter is
    transported back into the presynaptic terminal}
\end{keytermsframe}

% --- Key Term Card ---
\begin{frame}{Key Concept}
  \keyterm{Excitatory Post-Synaptic Potential (EPSP)}{%
    A transient depolarisation of the postsynaptic membrane caused by
    the opening of ligand-gated cation channels. EPSPs summate
    temporally and spatially to determine whether the postsynaptic
    neuron reaches firing threshold.}
\end{frame}

% --- Poll Question ---
\pollquestion{Which ion is most critical for triggering
  neurotransmitter vesicle fusion?}
  {Sodium (Na$^+$)}
  {Calcium (Ca$^{2+}$)}
  {Potassium (K$^+$)}
  {Chloride (Cl$^-$)}

% --- Worked Example ---
\begin{frame}{Calculating Reversal Potential}
  \begin{workedexample}
    \item Identify the ion: Na$^+$ with $[\text{Na}^+]_o = 145$\,mM
          and $[\text{Na}^+]_i = 12$\,mM
    \item Apply the Nernst equation:
          $E_{\text{Na}} = \frac{RT}{zF} \ln \frac{[\text{Na}^+]_o}{[\text{Na}^+]_i}$
    \item Substitute values at 37\textdegree C:
          $E_{\text{Na}} = \frac{61.5}{1} \log \frac{145}{12}$
    \item Calculate: $E_{\text{Na}} = 61.5 \times 1.08 \approx +66.5$\,mV
  \end{workedexample}
\end{frame}

% --- Exam Question ---
\begin{frame}{Sample Exam Question}
  \examquestion[15]{%
    Describe the sequence of events that occur during chemical synaptic
    transmission, from the arrival of an action potential at the
    presynaptic terminal to the generation of a postsynaptic potential.
    Include in your answer the role of voltage-gated calcium channels,
    SNARE proteins, and ligand-gated ion channels. Illustrate your
    answer with a clearly labelled diagram.}
\end{frame}

% --- Think--Pair--Share ---
\thinkpairshare{How might dysfunction at the synapse contribute to
  neurological or psychiatric disease? Give a specific example and
  explain the mechanism.}

% --- Recap ---
\begin{frame}{Session Summary}
  \begin{recap}
    \item Action potentials trigger Ca$^{2+}$ influx via voltage-gated
          channels at the presynaptic terminal
    \item SNARE proteins mediate vesicle fusion and neurotransmitter release
    \item Ionotropic receptors produce fast synaptic responses; metabotropic
          receptors produce slower, modulatory effects
    \item Synaptic transmission is terminated by reuptake, enzymatic
          degradation, or diffusion
  \end{recap}
\end{frame}

% --- Reading List ---
\begin{frame}{Preparation for Next Week}
  \begin{readinglist}
    \readingentry{Bear, M.~F., Connors, B.~W., \& Paradiso, M.~A.}
      {2020}{Neuroscience: Exploring the Brain}{Jones \& Bartlett, Ch.~5--6}
    \readingentry{Kandel, E.~R.~et~al.}{2021}
      {Principles of Neural Science}{McGraw-Hill, Ch.~10--12}
  \end{readinglist}
\end{frame}

\conclusionslide{See You Next Week}{james.obrien@manchester.ac.uk}

\end{document}
